\section[Cote 77 : le discriminant universel modulo 2 et modulo 4, la symétrie canonique]{Cote 77 : le discriminant universel modulo $2$ et modulo $4$, la symétrie canonique\protect\verifie{Folder77}}\label{sec:77}

La cote 77, \emph{Quadriques affines, extensions quadratiques : notes manuscrites (s.d.)}, compte cent cinquante-sept pages, que l'inventaire place après 1977. La partie IV, dont la couverture de la page~64 porte \q{discriminant des formes quadratiques}, prend la forme quadratique universelle $\sum_i Y_i x_i^2 + \sum_{i<j} Y_{ij} x_i x_j$ sur l'anneau $\mathbf Z[Y]$ et le déterminant $\Delta$ de sa forme polaire, matrice qui porte $2Y_i$ sur la diagonale et $Y_{ij}$ ailleurs. Les pages~65 à~67 développent $\Delta$ en une somme sur des graphes : sommets isolés pour les points fixes, arêtes doubles pour les transpositions, polygones pour les autres cycles. La page~68 réduit cette somme modulo~$2$ : il n'y reste que les involutions sans point fixe, et la page conclut \q{On voit donc que $\Delta \bmod 2 = 0$ si $\operatorname{card} I$ est impair} (\lot{77}{4}{68}). En rang pair $m = 2n$, elle pose $B$, somme sur les couplages parfaits de $I$ des produits des $Y_{ij}$, cherche $\varepsilon = \pm 1$ tel que $B^2 - \varepsilon\Delta \equiv 0 \pmod 4$, et \q{prévoit} $\varepsilon = (-1)^n$. Les pages~69 et~71 comparent $B^2$ et $\Delta$ modulo~$4$ par les paires de couplages, avec une fausse alarme à l'encre brune corrigée deux lignes plus bas. La page~87, barrée, contrôle les coefficients du développement en rang~$4$ (\lot{77}{5}{87}). La page~127 prend le discriminant divisé $\delta' = \delta/2$ du rang impair comme critère de régularité (\lot{77}{7}{127}). La partie V, aux pages~117 à~125, étudie les quadriques affines \q{birégulières}, lisses et lisses à l'infini : elle les ramène à un triplet $(T, q, a)$ avec $q(a) = 1$, pose $\pi = \varphi(a, \cdot)$ et, page~123, définit la symétrie canonique $\sigma = \mathrm{id} - \pi \otimes a$ ; une note en marge affirme qu'elle est la seule réflexion qui fixe l'hyperplan à l'infini et laisse la quadrique invariante (\lot{77}{7}{123}). Le feuillet découpé de la page~124 traite l'exemple $T = M_2$, $q = \det$, $a = 1$.

\noindent\emph{Conventions.} Pour $n \geq 0$, les indéterminées sont indexées par les parties $J \subset \{1, \dots, n\}$ de cardinal $1$ ou $2$, notées $Y_i$ et $Y_{ij}$ ; on note $\mathbf Z[Y]$ l'anneau de polynômes correspondant. La matrice polaire universelle $P_n$ a pour coefficients $2Y_i$ en $(i, i)$ et $Y_{ij}$ en $(i, j)$ et $(j, i)$ pour $i \neq j$. Sur un anneau commutatif $R$, une famille $y = (y_i, y_{ij})$ d'éléments de $R$ définit la forme $\sum y_i x_i^2 + \sum_{i<j} y_{ij} x_i x_j$, de matrice polaire $P(y)$, image de $P_n$ par le morphisme $Y \mapsto y$.

\begin{enonce}[Page 68, rang impair\piste{77-odd-rank-half-discriminant}]
Si $n$ est impair, $\det P_n$ est divisible par $2$ dans $\mathbf Z[Y]$. Le quotient $\delta'_n = \det P_n / 2$ est l'unique polynôme tel que $\det P_n = 2\delta'_n$, et pour tout anneau commutatif $R$ et toute famille $y$ d'éléments de $R$, $\det P(y) = 2\,\delta'_n(y)$.
\end{enonce}

\begin{proof}
Soit $K_n$ la matrice de $\mathbf Z[Y]$ qui porte $Y_{ij}$ en $(i, j)$ pour $i < j$, $-Y_{ij}$ en $(i, j)$ pour $i > j$ et $0$ sur la diagonale. Elle est antisymétrique : ${}^tK_n = -K_n$. On a donc $\det K_n = \det {}^tK_n = \det(-K_n) = (-1)^n \det K_n = -\det K_n$, d'où $2 \det K_n = 0$ ; comme $\mathbf Z[Y]$ est intègre de caractéristique nulle, $\det K_n = 0$.

Soit $\rho : \mathbf Z[Y] \to \mathbf F_2[Y]$ la réduction des coefficients modulo~$2$. Les matrices $P_n$ et $K_n$ ont même image par $\rho$ : sur la diagonale, $\rho(2Y_i) = 0$ ; au-dessus, les coefficients sont égaux ; au-dessous, $\rho(-Y_{ij}) = \rho(Y_{ij})$ en caractéristique~$2$. Le déterminant commutant aux morphismes d'anneaux, $\rho(\det P_n) = \det \rho(P_n) = \det \rho(K_n) = \rho(\det K_n) = 0$. Le noyau de $\rho$ est formé des polynômes dont tous les coefficients sont pairs, c'est-à-dire de l'idéal engendré par~$2$ ; donc $2$ divise $\det P_n$.

L'unicité du quotient vient de ce que $2$ n'est pas diviseur de zéro dans $\mathbf Z[Y]$. Pour $R$ et $y$ donnés, le morphisme $Y \mapsto y$ envoie $P_n$ sur $P(y)$ et commute au déterminant ; il envoie $2\delta'_n$ sur $2\,\delta'_n(y)$.
\end{proof}

\begin{remarque}
L'énoncé tient pour tout $n$ impair. La démonstration ne suit pas le développement en graphes de la page : elle remplace, modulo~$2$, la matrice symétrique par la matrice antisymétrique $K_n$, dont le déterminant est nul sur $\mathbf Z[Y]$ en taille impaire. Sur $\mathbf F_2[Y]$ l'argument $\det = -\det$ ne donne rien, et c'est pourquoi il faut relever $\rho(P_n)$ en une matrice antisymétrique sur $\mathbf Z[Y]$. La seconde partie de la piste, \q{$q$ est régulière si et seulement si $\delta'(q)$ est inversible}, est la définition de la régularité que la lecture adopte (page~127). Les équivalences de la page~127 avec la lissité de $V(q)$ et l'isomorphie locale à la forme standard ne sont pas formalisées.
\end{remarque}

\begin{enonce}[Pages 68 à 71 et 87, rangs $2$ et $4$\piste{77-universal-even-rank-discriminant-mod-4}]
Soit $R$ un anneau commutatif.
\begin{enumerate}
\item (Rang $2$, $n = 1$.) Pour $a_1, a_2, b_{12} \in R$, $-\det\begin{pmatrix} 2a_1 & b_{12} \\ b_{12} & 2a_2\end{pmatrix} = B^2 - 4C$ avec $B = b_{12}$ et $C = a_1a_2$.
\item (Rang $4$, $n = 2$.) Soit $P$ la matrice polaire de la forme $\sum_{i \leq 4} a_i x_i^2 + \sum_{i<j} b_{ij} x_i x_j$ et $B = b_{12}b_{34} + b_{13}b_{24} + b_{14}b_{23}$, somme sur les trois couplages parfaits de $\{1, 2, 3, 4\}$. Alors $\det P = B^2 - 4C$ avec
\begin{multline*}
C = -4a_1a_2a_3a_4 + a_1a_2b_{34}^2 + a_1a_3b_{24}^2 + a_1a_4b_{23}^2 + a_2a_3b_{14}^2 + a_2a_4b_{13}^2 + a_3a_4b_{12}^2 \\
- a_1b_{23}b_{24}b_{34} - a_2b_{13}b_{14}b_{34} - a_3b_{12}b_{14}b_{24} - a_4b_{12}b_{13}b_{23} \\
+ b_{12}b_{13}b_{24}b_{34} + b_{12}b_{14}b_{23}b_{34} + b_{13}b_{14}b_{23}b_{24}.
\end{multline*}
\item Le signe $\varepsilon = -(-1)^n$ ne convient dans aucun des deux rangs : sur $\mathbf Z$, pour $a_1 = a_2 = 0$ et $b_{12} = 1$, $B^2 - \det P = 2$ ; pour $b_{12} = b_{34} = 1$ et les autres coefficients nuls, $B^2 + \det P = 2$. Aucun des deux n'est divisible par~$4$.
\end{enumerate}
\end{enonce}

\begin{proof}
En rang~$2$, $\det = 4a_1a_2 - b_{12}^2$. En rang~$4$, on développe $\det P$ selon la première ligne, puis chaque mineur d'ordre~$3$ par la règle de Sarrus. On obtient
\begin{multline*}
\det P = 16a_1a_2a_3a_4 - 4\,(a_1a_2b_{34}^2 + a_1a_3b_{24}^2 + a_1a_4b_{23}^2 + a_2a_3b_{14}^2 + a_2a_4b_{13}^2 + a_3a_4b_{12}^2) \\
+ 4\,(a_1b_{23}b_{24}b_{34} + a_2b_{13}b_{14}b_{34} + a_3b_{12}b_{14}b_{24} + a_4b_{12}b_{13}b_{23}) + b_{12}^2b_{34}^2 + b_{13}^2b_{24}^2 + b_{14}^2b_{23}^2 \\
- 2\,(b_{12}b_{13}b_{24}b_{34} + b_{12}b_{14}b_{23}b_{34} + b_{13}b_{14}b_{23}b_{24}).
\end{multline*}
Ce sont les coefficients de la page~87 : $2^4$ pour les quatre sommets isolés, $-2^2$ pour une arête double et deux sommets isolés, $2^2$ pour un triangle et un sommet isolé, $1$ pour deux arêtes doubles, $-2$ pour un carré. Le carré $B^2$ a pour termes carrés les $b_{12}^2b_{34}^2$, $b_{13}^2b_{24}^2$, $b_{14}^2b_{23}^2$ et pour termes croisés $2\,(b_{12}b_{13}b_{24}b_{34} + b_{12}b_{14}b_{23}b_{34} + b_{13}b_{14}b_{23}b_{24})$ ; la différence $B^2 - \det P$ est donc $4C$, terme à terme. Les deux contre-exemples se calculent sur ces formules : en rang~$2$, $B^2 - \det P = 1 - (-1)$ ; en rang~$4$, $\det P = b_{12}^2b_{34}^2 = 1$ et $B^2 + \det P = 2$.
\end{proof}

\begin{remarque}
L'identité est démontrée sur tout anneau commutatif, donc en particulier sur $\mathbf Z[Y]$, avec un $C$ écrit en entier. Elle confirme en rangs $2$ et $4$ le signe $(-1)^n$ que la page \q{prévoit}, et le choix opposé échoue dans ces deux rangs, ce que la lecture avait contrôlé par machine pour $m = 2, 4, 6$. En rang~$4$, la vérification confirme les coefficients et les nombres de graphes de la page~87. Le rang pair quelconque fait l'objet de l'énoncé suivant, sans expression de $C$.
\end{remarque}

\begin{enonce}[Pages 68 à 71, rang pair quelconque\piste{77-universal-even-rank-discriminant-mod-4}]
Soient $n \geq 0$ et $B_{2n} = \sum_{M} \prod_{\{i, j\} \in M} Y_{ij}$, la somme portant sur les couplages parfaits $M$ de $\{1, \dots, 2n\}$. Alors $4$ divise $(-1)^n \det P_{2n} - B_{2n}^2$ dans $\mathbf Z[Y]$. Il existe donc $C_{2n} \in \mathbf Z[Y]$ tel que, pour tout anneau commutatif $R$ et toute famille $y$ d'éléments de $R$, $(-1)^n \det P(y) = B_{2n}(y)^2 - 4\,C_{2n}(y)$.
\end{enonce}

\begin{proof}
Soit $U$ la matrice qui porte $Y_i$ en $(i, i)$, $Y_{ij}$ en $(i, j)$ pour $i < j$, et $0$ sous la diagonale. Alors $P_{2n} = U + {}^tU$, et la matrice antisymétrique $K_{2n}$ de l'énoncé du rang impair est $U - {}^tU = P_{2n} - 2\,{}^tU$. On écrit $P$, $K$ et $B$ pour $P_{2n}$, $K_{2n}$ et $B_{2n}$.

\emph{Premier pas : $\det K \equiv (1 - 2n) \det P \pmod 4$.} Soit $f(t) = \det(P - t\,{}^tU) \in \mathbf Z[Y][t]$, que l'on écrit $f(t) = f_0 + f_1 t + t^2 g(t)$. On a $f_0 = f(0) = \det P$ et $\det K = f(2) = \det P + 2f_1 + 4g(2)$ ; il reste à voir que $2f_1 = -2n \det P$. Le déterminant $\det P$ n'est pas nul : pour $Y_i = 1$ et $Y_{ij} = 0$ il vaut $2^{2n}$. La matrice $P$ est donc inversible sur le corps des fractions $L$ de $\mathbf Z[Y]$, et sur $L$, $f(t) = \det P \cdot \det(1 - t\,P^{-1}\,{}^tU)$, dont le coefficient de $t$ est $-\det P \cdot \operatorname{tr}(P^{-1}\,{}^tU)$. Comme $P$ est symétrique, $P^{-1}$ l'est aussi, et $\operatorname{tr}(P^{-1}\,{}^tU) = \operatorname{tr}\bigl({}^t(P^{-1}\,{}^tU)\bigr) = \operatorname{tr}(U P^{-1}) = \operatorname{tr}(P^{-1} U)$. La somme de ces deux traces est $\operatorname{tr}(P^{-1}(U + {}^tU)) = \operatorname{tr}(1) = 2n$. Donc $2f_1 = -2n \det P$ dans $L$, et dans $\mathbf Z[Y]$, qui s'injecte dans $L$.

\emph{Deuxième pas : $\det K$ est un carré dans $\mathbf Z[Y]$.} Sur un corps $k$ de caractéristique différente de~$2$, une matrice antisymétrique $A$ de taille paire a pour déterminant un carré. On raisonne par récurrence sur la taille ; la diagonale de $A$ est nulle. Si la première ligne est nulle, $\det A = 0$. Sinon, soit $c = A_{1j} \neq 0$, avec $j \neq 1$ ; la permutation qui échange $2$ et $j$, appliquée aux lignes et aux colonnes, ne change pas le déterminant et amène $c$ en position $(1, 2)$. On découpe alors $A$ en blocs $\begin{pmatrix} E & F \\ G & D \end{pmatrix}$, avec $E = \begin{pmatrix} 0 & c \\ -c & 0 \end{pmatrix}$ inversible de déterminant $c^2$ ; on a $\det A = c^2 \det S$ avec $S = D - G E^{-1} F$. Comme $G = -{}^tF$, ${}^tD = -D$ et ${}^t(E^{-1}) = -E^{-1}$, la matrice $S$ est antisymétrique, de taille deux de moins. Par récurrence $\det S = s^2$, d'où $\det A = (cs)^2$. Appliqué à l'image de $K$ sur $L$, cela donne $\det K = s^2$ avec $s \in L$. L'élément $s$ est racine du polynôme unitaire $T^2 - \det K$ de $\mathbf Z[Y][T]$ ; l'anneau $\mathbf Z[Y]$ est factoriel, donc intégralement clos, et $s$ est l'image d'un $r \in \mathbf Z[Y]$, avec $\det K = r^2$.

\emph{Troisième pas : réduction modulo~$2$.} Soit $M$ une matrice symétrique à diagonale nulle sur un anneau de caractéristique~$2$. Dans $\det M = \sum_\sigma \varepsilon(\sigma) \prod_i M_{\sigma(i) i}$, les signes valent~$1$, et $\sigma$ et $\sigma^{-1}$ donnent le même produit, par symétrie de $M$. Les termes des $\sigma$ distinctes de $\sigma^{-1}$ s'annulent donc deux à deux ; ceux des involutions qui ont un point fixe sont nuls, car la diagonale est nulle. Il reste les involutions sans point fixe, c'est-à-dire les couplages parfaits, et pour un couplage le produit est le carré du produit des $M_{ij}$ sur ses arêtes. En caractéristique~$2$ une somme de carrés est le carré de la somme, d'où $\det M = \bigl(\sum_{M'} \prod_{\{i, j\} \in M'} M_{ij}\bigr)^2$. La matrice $\rho(K)$ est symétrique, car $-1 = 1$ dans $\mathbf F_2[Y]$, à diagonale nulle, et ses coefficients au-dessus de la diagonale sont les $\rho(Y_{ij})$ ; on obtient $\rho(r)^2 = \rho(\det K) = \rho(B)^2$. En caractéristique~$2$, $\rho(r - B)^2 = \rho(r)^2 - \rho(B)^2 = 0$, et comme $\mathbf F_2[Y]$ est intègre, $\rho(r - B) = 0$. Comme dans l'énoncé du rang impair, $r = B + 2h$ avec $h \in \mathbf Z[Y]$, et $\det K = r^2 = B^2 + 4(Bh + h^2)$.

\emph{Conclusion.} Par récurrence sur $n$, $(-1)^n \equiv 1 - 2n \pmod 4$. Donc $(-1)^n \det P \equiv (1 - 2n) \det P \equiv \det K \equiv B^2 \pmod 4$. Si $(-1)^n \det P - B^2 = 4Q$, on prend $C_{2n} = -Q$ ; le morphisme $Y \mapsto y$ envoie $P_{2n}$ sur $P(y)$ et commute au déterminant.
\end{proof}

\begin{remarque}
L'énoncé tient en tout rang pair, avec le signe $(-1)^n$ que la page \q{prévoit}. La démonstration ne suit pas les pages~69 et~71, qui comparent $B^2$ et $\Delta$ par les paires de couplages. Elle passe par la matrice antisymétrique $K = U - {}^tU$. Le signe $(-1)^n$ vient du terme du premier ordre, $\det K \equiv (1 - 2n) \det P \pmod 4$ ; la congruence $\det K \equiv B^2$ vient de ce que $\det K$ est le carré d'un polynôme $r$ qui se réduit sur $B$ modulo~$2$. Ce polynôme $r$ est, au signe près, le pfaffien de $K$, qui manque à mathlib ; il n'est pas construit ici, son existence découle de la clôture intégrale de $\mathbf Z[Y]$. Le polynôme $C_{2n}$ n'est écrit qu'en rangs $2$ et~$4$.
\end{remarque}

\begin{enonce}[Pages 117 à 124, la symétrie canonique\piste{77-biregular-affine-quadric-symmetry}]
Soient $R$ un anneau commutatif, $T$ un $R$-module, $q$ une forme quadratique sur $T$, $\varphi(x, y) = q(x + y) - q(x) - q(y)$ sa forme polaire, $a \in T$ avec $q(a) = 1$, $\pi = \varphi(a, \cdot)$, $V = \ker \pi$, $E = \pi^{-1}(1)$ et $\sigma(x) = x - \pi(x)\,a$.
\begin{enumerate}
\item On a $\pi(a) = 2$, $q \circ \sigma = q$ et $\sigma^2 = \mathrm{id}$, et $\sigma$ fixe $V$ point par point.
\item Supposons $\varphi$ non dégénérée, au sens où $\varphi(y, \cdot) = 0$ entraîne $y = 0$, et $\pi$ surjective. Une application linéaire $u : T \to T$ qui fixe $V$ point par point et vérifie $q \circ u = q$ est de la forme $u(x) = x - e\,\pi(x)\,a$ avec $e \in R$ idempotent. Réciproquement, pour tout idempotent $e$, cette application préserve $q$ et fixe $V$. Si $0$ et $1$ sont les seuls idempotents de $R$, alors $u = \mathrm{id}$ ou $u = \sigma$.
\item Si $2$ est inversible dans $R$, alors $\pi(a/2) = 1$, et pour $x \in E$, $-\sigma(x) \in E$ et $-\sigma(x) - a/2 = -(x - a/2)$. Sur $E$, la transformation projective définie par $\sigma$ est la symétrie de centre $a/2$.
\item Si $2 = 0$ dans $R$, alors $\pi(a) = 0$, donc $a \in V$, $\pi \circ \sigma = \pi$, et $\sigma(x) = x - a$ pour $x \in E$ : $\sigma = \mathrm{id} - \pi \otimes a$ est une transvection.
\item (Page~124.) Pour $T = M_2(R)$, $q = \det$ et $a = 1$, $\varphi(x, 1) = \det(x + 1) - \det x - 1 = \operatorname{tr} x$.
\end{enumerate}
\end{enonce}

\begin{proof}
On a $\varphi(a, a) = 2q(a) = 2$. Pour $t \in R$, $q(x + ta) = q(x) + t\,\varphi(a, x) + t^2 q(a)$ ; avec $t = -\pi(x)$, $q(\sigma x) = q(x) - \pi(x)^2 + \pi(x)^2 = q(x)$. Ensuite $\pi(\sigma x) = \pi(x) - 2\pi(x) = -\pi(x)$, d'où $\sigma(\sigma x) = \sigma x + \pi(x)\,a = x$. Si $\pi(v) = 0$, $\sigma v = v$.

Pour le point~2, soient $b$ tel que $\pi(b) = 1$ et $c = u(b) - b$. Pour tout $x$, $x - \pi(x) b$ est dans $V$, donc $u(x) = x + \pi(x)\,c$. Alors $q(u x) = q(x) + \pi(x) f(x)$ avec $f(x) = \pi(x) q(c) + \varphi(x, c)$, forme linéaire, et l'hypothèse donne $\pi(x) f(x) = 0$ pour tout $x$. En $x = b$, $f(b) = 0$. En $x + b$, $0 = (\pi(x) + 1)\bigl(f(x) + f(b)\bigr) = \pi(x) f(x) + f(x) = f(x)$. Donc $\varphi(c + q(c)\,a, \cdot) = f = 0$, et $c = -q(c)\,a$ par non-dégénérescence. Posons $e = q(c)$ ; alors $e = q(-e a) = e^2$, et $u(x) = x - e\,\pi(x)\,a$. Réciproquement, $q(x - e\,\pi(x)\,a) = q(x) - e\,\pi(x)^2 + e^2 \pi(x)^2 = q(x)$. Si $e = 0$, $u = \mathrm{id}$ ; si $e = 1$, $u = \sigma$.

Pour le point~3, $\pi(a/2) = \pi(a)/2 = 1$ ; si $\pi(x) = 1$, $\sigma(x) = x - a$, $\pi(-\sigma x) = -(1 - 2) = 1$ et $-\sigma(x) - a/2 = a/2 - x$. Pour le point~4, $\pi(a) = 2 = 0$, $\pi(\sigma x) = \pi(x) - \pi(x)\,\pi(a) = \pi(x)$, et $\sigma(x) = x - a$ si $\pi(x) = 1$. Pour le point~5, $\det(x + 1) = \det x + \operatorname{tr} x + 1$ pour une matrice $2 \times 2$.
\end{proof}

\begin{remarque}
La lecture adopte la normalisation $\pi = \varphi(a, \cdot)$, que la page~117 biffe. C'est elle qui donne $\pi(a) = 2$ et $\sigma^2 = \mathrm{id}$ ; la page justifie (21) par \q{$\pi(a) = 1$}, qui donnerait $\sigma^2 = \sigma$. L'unicité utilise deux hypothèses que la page tire de sa situation : $\varphi$ non dégénérée (\q{comme $\psi$ est iso}) et $\pi$ surjective, ce que donne la structure d'extension (1). Elle utilise aussi l'absence d'idempotent non trivial. La page conclut de $\lambda(\lambda + 1) = 0$ que $\lambda = -1$ hors de l'identité ; sur un anneau quelconque les solutions sont les $\lambda = -e$, $e$ idempotent. Sur $R = R_1 \times R_2$ et pour $e = (1, 0)$, le point~2 donne une isométrie qui fixe $V$ et qui vaut $\sigma$ au-dessus de la première composante et l'identité au-dessus de la seconde. La Proposition de la marge vaut donc sur une base connexe, ou fibre à fibre. L'énoncé formalisé traite l'invariance $q \circ u = q$ ; l'invariance à une unité près, $q \circ u = \lambda q$, dont la page déduit $\lambda = 1$ en rang au moins~$3$, n'est pas formalisée. Sur la page~124 la polaire de $\det$ est écrite avec le signe opposé, d'où $\pi = -\operatorname{tr}$ et la remarque \q{attention au signe !} ; avec la convention $\varphi(x, y) = q(x + y) - q(x) - q(y)$, $\pi$ est la trace et $\pi(1) = 2$.
\end{remarque}

Ne sont pas démontrés ici : une expression de $C$ en rang pair $2n \geq 6$ ; le développement de $\Delta$ en graphes des pages~65 à~67 ; les équivalences de la page~127 entre l'inversibilité de $\delta'(q)$, la lissité de $V(q)$ et l'isomorphie locale à la forme standard ; l'invariance de $q$ à une unité près et $\lambda = 1$ en rang au moins~$3$ (page~123) ; l'équivalence du Scholie (page~120) entre les $(P, Q, H)$ et les $(T, q, a)$ ; la symétrie (22) du rang impair et la structure $G \simeq G^0 \times \mathbf Z/2$ annoncée page~125.
